\section{Operation instructions}

The SimpleSBC platform is intended to provide a low-cost embedded Linux 
development environment for research, education, and custom instrumentation 
applications. Once assembled and programmed, the board may be operated either 
as a standalone Linux system or as the processing component of a larger 
experimental platform through the use of the expansion and peripheral 
interfaces.

\subsection{Preparing the system}

Prior to operation, verify that the board has been assembled correctly 
according to the instructions provided in the Build Instructions section. 
Confirm that a valid Linux image has been written to a microSD card using 
the Buildroot configuration supplied in the repository.

Insert the microSD card into the microSD socket and ensure that any required 
peripherals, such as displays, USB devices, serial connections, or daughter 
boards, are connected before power is applied.

\begin{figure}[h]
\centering
\includegraphics[width=0.8\textwidth]{Images/board_overview.png}
\caption{Overview of the assembled SimpleSBC platform showing major connectors and interfaces.}
\label{fig:board_overview}
\end{figure}

Figure \ref{fig:board_overview} should identify the microSD connector, USB-C power input, 
serial interface, user buttons, and any expansion connectors available on the board.

\subsection{Powering the board}

Connect a suitable USB power source to the USB-C connector. 
During development, it is recommended to use a current-limited bench supply or 
a USB power meter to monitor power consumption.

After power is applied, the board will begin the boot sequence automatically. 

Typical boot times are less than one minute depending on the software image and attached peripherals.
If a serial terminal is connected through the RP2040 USB-to-UART bridge, 
Linux boot messages should appear shortly after power is applied. 
Successful completion of the boot sequence is indicated by the appearance of a 
login prompt or application-specific startup messages.

\subsection{Serial console access}

During development and debugging, interaction with the system is typically 
performed through the serial console.

\begin{enumerate}
\item Connect the USB cable to the host computer.
\item Identify the serial device created by the RP2040 USB-to-UART bridge.
\item Open a terminal emulator such as PuTTY, Tera Term, minicom, screen, or picocom.
\item Configure the serial connection according to the settings specified by the software image.
\item Reset or power cycle the board if boot messages are not visible.
\end{enumerate}

The serial console provides access to Linux boot messages, diagnostic output, 
and command-line administration of the system.

\subsection{Running software}

After the system has booted successfully, user applications may be executed 
directly from the Linux command line or configured to launch automatically 
during startup.

The Buildroot image included with this work contains demonstration scripts 
and example applications that may be used to validate correct operation of 
the processor, storage system, display interface, and software environment. 
Researchers may replace these applications with custom software appropriate 
to their own experimental requirements.

The board was intentionally designed to support customization. Additional 
peripherals may be connected through the available expansion interfaces and 
controlled through Linux device drivers, user-space applications, or custom 
kernel modifications.

\subsection{Using custom daughter boards}

A significant portion of the PCB area is reserved for daughter-board integration. 
Researchers wishing to adapt the platform for specific experiments may install 
custom hardware in this area and interface with the Linux system through GPIO, 
UART, SPI, I\textsuperscript{2}C, PWM, or other available interfaces.

When integrating custom hardware:

\begin{enumerate}
\item Verify voltage compatibility with the board.
\item Confirm pin assignments using the schematic and device tree files.
\item Test new peripherals individually before integrating them into larger systems.
\item Update the Linux device tree if additional hardware resources are required.
\end{enumerate}

\subsection{Shutdown procedure}

To prevent filesystem corruption, the operating system should be shut down 
gracefully before removing power.

From the Linux command line:

\begin{verbatim}
shutdown -h now
\end{verbatim}

or

\begin{verbatim}
poweroff
\end{verbatim}

Wait until all filesystem activity has stopped before disconnecting power 
or removing the microSD card.

\subsection{Troubleshooting}

If the board does not boot:

\begin{enumerate}
\item Verify that the microSD card contains a valid software image.
\item Confirm that the power supply is capable of providing sufficient current.
\item Check for visible assembly defects or damaged components.
\item Monitor the serial console for boot messages.
\item Verify that the correct device tree and Buildroot configuration were 
used during image generation.
\end{enumerate}

If the display or peripheral hardware does not function correctly:

\begin{enumerate}
\item Verify connector orientation and cable installation.
\item Confirm that the required kernel patches and device tree configuration 
are present.
\item Check Linux system logs for driver initialization messages.
\item Verify peripheral functionality independently where possible.
\end{enumerate}

\subsection{Safety considerations}

The platform operates at low voltages and is intended for laboratory, 
educational, and research use. Nevertheless, users should avoid modifying 
the hardware while power is applied and should verify that all external 
peripherals operate within the supported voltage ranges.

Improper wiring of expansion hardware may damage the processor, peripheral 
devices, or power-management circuitry. During initial development, current 
limiting is recommended whenever new hardware is connected.

The board should not be exposed to moisture, conductive debris, or excessive 
mechanical stress. If any component becomes unusually hot during operation, 
disconnect power immediately and investigate the cause before further use.
